\chapter[Sermon 5. The Beginning of the Work Is Made, and a Most Goodly Description to Us Exhibited of Christ King and Bishop in Glory, and Nevertheless Working in the Church.]{The Beginning of the Work Is Made, and a Most Goodly Description to Us Exhibited of Christ King and Bishop in Glory, and Nevertheless Working in the Church.}
\label{sermon:005}
\markright{Sermon 5. The Beginning of the Work Is Made, and a Most Goodly ...}

And I turned myself to look at the voice that spoke with me. And when I turned, I saw seven golden lampstands. And in the midst of the seven lampstands, I saw one like the Son of Man, clothed in a linen garment that reached to the ground, and girded about the chest with a golden belt. His head and his hair were white, like white wool and snow. And his eyes were like a flame of fire, and his feet were like burnished bronze, as if they had been heated in a furnace. And his voice was like the sound of many waters. And he had in his right hand seven stars, and out of his mouth went out a sharp two-edged sword, and his face shone as the sun in its strength.

Such things as have been treated of hitherto in this book serve in place of a prologue or preface, as they term it. Now at last the matter itself will be presented to us. Therefore, what follows is the second part of this book, which extends to the fourth chapter. In this, Christ is described to us with his universal church. For first, we are shown the most sacred image of Christ our Lord, teaching what he is on the right hand of his Father in glory, and how he, sitting on the right hand of his Father, continually works in his church, never absent, present always. Of what sort the church is here on earth, is figured in those seven congregations. Here, therefore, are shown the excellent gifts of churches, and again, shameful errors: How the Lord Christ confirms those who are sliding and ready to fall, establishes those that stand, strengthens the weakhearted, restrains the foolish and bold, and preserves things that are corrupt. Finally, how faithful ministers of the church must work and travel with the people committed to their care. For here is exceedingly taught what is the repairing and preservation of churches. Where also a brief summary of the whole ecclesiastical doctrine, brought in to an abridgement, shall be set before us. For here is repeated from heaven concerning Christ in glory, the doctrine of true religion, which he had set forth more plentifully when he was yet here on earth; and here it is most aptly applied to churches, after consideration of the same.

\index{Arethas of Caesarea}And in most goodly order the words are knit together (as likewise the whole book is written with plain words and hanging right well together; they are deceived who think it to be loose, unbound writings). John heard a voice behind him crying: whereupon he turned to see the voice speaking, that is to wit, him who spake. For Arethas also admonishes that there is a trope. For no man sees, but hears the voice. After turning to see, he saw a figure of Christ our Savior. Therefore, when the Lord speaks, let us turn also with all our heart, that we may likewise deserve to see the mysteries of the kingdom of God, for he gladly reveals himself to those who turn and desire heavenly things. And from those who neglect the mysteries of the kingdom of God, all things of salvation are hidden.

Furthermore, John depicts to us the image of Christ, our universal king and chief bishop, seated in glory: and within this description are contained the principal matters concerning Christ. For such a glimpse of Christ is here given us as our frail flesh may perceive in this world. But we shall see him ultimately in the world to come as he is, in the fullness of his majesty, wherein shall be joy and everlasting life; but this in this corrupt world is yet granted to no man.

So much, therefore, is permitted to us who live yet in this world and are still visible, as is profitable, and as our weakness may comprehend. But this is not small or insignificant; it is great, abundant, and full of spiritual joy, if we behold these mysteries of God with a faithful eye and a mind eager for godly matters. And there is no doubt that these are certain and true things that are revealed to us here. For they are revealed by the very Son of God. Let us not wish to see more, or desire greater things than these are, but take pleasure in those which Christ has granted us. And let us know for certain that a wonderful benefit of God is given to us in this vision. For who would not desire to see Christ in glory, sitting on the right hand of the Father? Who does not desire to know what our Savior does in heaven? And how, being in heaven, he is nevertheless present with his church on earth? But this sacred and holy image instructs in all these points most fully the faithful of Christ. However, this image of Christ is not to be depicted with colors, since colors cannot attain to its majesty, but with the ecclesiastical doctrine, which has the promise of the Spirit of Christ, and is therefore more evident and fitting for its true expression. Let us also print this image, not upon any dead surface with colors that will perish and fade, but in our hearts through the living Spirit of God, which may also keep it in our minds, never to be erased. And the things spoken in the second and third chapters of this book are derived from this description of Christ, that the majesty of the thing might invite us to singular diligence. The matter is very plain.

First we are taught who it is, whose image is to be exhibited to us: not the Son of Man himself in his own substance, but like the Son of Man. The Son of Man, according to the Gospel, is called Christ himself, very God. Here he did not show himself to be seen as John in his own substance, but in the form of an angel, that represents Christ. This thing is found more than once in the book. We shall therefore refer all these things unto Christ, not to the angel, which is the minister of Christ in the mystery. And we shall see Christ in his own substance, what time our base body shall depart from hence, and being raised from the dead shall be glorified. In the meantime, the soul from the death of the body till it rises again, shall clearly have the fruition of the sight of Christ: wherein, as I said before, shall be the chief joy and felicity. We shall now therefore see Christ as it were in a glass, and so much as shall suffice us. The Lord open to us the eyes of our mind.

\index{Augustine, Saint}He also tells where he saw Christ, in the midst of seven candlesticks. Soon we will see that the candlesticks must be understood as representing the churches. Christ is then in the midst of the church. He sits truly on the right hand of the Father, and according to the nature of his divine body, he is only in one place, and no more. As Saint Augustine declares abundantly in the fifty-seventh Epistle to Donatus. Yet, because he is also very God, he is likewise in the midst of the church, as he promised in the Gospel: “Wherever two or three are gathered in my name, I am in the midst of them.” And again: “Behold, I am with you until the end of the world.” Therefore, by his power, Christ remains and works in the present church, and is not absent. (Leave therefore to inquire, what Christ does on the right hand of his Father, whether he sits continually?) And he is truly in the midst of the churches, fixed to a place, but shows himself equally to all, as helpful. For he accepts no persons, nor sleeps. He is not painted, he is not idle, paying no regard to matters of the church, but is chiefly and only attentive to the salvation of the same. Such a one he promised himself to be in chapters 14, 15, and 16 of John. And saying that Christ is in the midst of the church, what Vicar shall he have? Shall he have that enemy which is directly against him? For a Vicar is in the place of one who is absent: but Christ is in the middle of the church, present, not absent.

\index{Jerome, Saint}\index{Justification}\index{Zechariah (Book of)}\index{Hebrews, Epistle to the}Christ is described most plentifully here, and many things are ascribed to him. It is declared in what manner Christ is in the midst of the church. First, it is shown what garment he has on: namely, both priestly and princely. By this, it is figured what manner of one Christ is in heaven and on earth: namely, bishop and king, intercessor, mediator, and sacrifice, a most perfect sanctification and justification, a redeemer and deliverer of the faithful to his father, evermore working the salvation of his faithful. As Paul teaches, Romans 8, and Hebrews 7. A \textit{poderes} [?] is found among the apparel of Aaron, and it is a priestly garment. Jerome writes to Fabiola concerning the priestly garment. The second vesture of linen is a coat down to the foot, of double linen. Josephus calls it \textit{bessina}. It is called in Hebrew \textit{ketheneth}, in Greek [illegible]. This clings closely to the body, and is so narrow and straight-cut that there is no wrinkle at all in the garment, and came down to the legs. This was truly white and clean. For the Lord Christ is an undefiled priest, as in Hebrews 7. Neither does he wear again a foul vesture, as he did in Zechariah 3, nor a purple, as in John 19. But a bright one, as he who has obtained a name above all names. But his girdle or belt is worn by soldiers and triumphant persons, and it signifies in Christ royal dignity. For Christ is king, deliverer, and redeemer of the faithful. His victory is ours. He has overcome Satan, Hell, sin, and death.

But the belt or girdle of Christ is not set in the customary place, namely, about the loins. For as Arethas has also admonished, there are no desires to be restrained in Christ. Therefore, he is not girded after the manner of sinners, but about the breasts: so that we should understand by the girding that he is king of kings, devoid of all affections; most righteous and holy in judgments and government. But yet, in the meantime, he is furnished for the defense of his church, as we have read it written in the 9th Psalm. The Lord has put on strength and girded himself, and so forth. Christ might seem to have girded himself not after the manner that priests or kings use, for that he has obtained a more excellent priesthood and kingdom, enduring forever. To accomplish these things, it behooved him to use a temple and palace, not transitory, but heaven itself. Hebrews 6 and 9. Yet, in the meantime, the effect penetrates into the church itself, so that he may be present in the church also.

\index{Papacy}\index{Daniel (Book of)}\index{Rome}But the head of Christ appears white, and his hair white, like the finest wool and the whitest snow. Such a head is also ascribed to the Father of our Lord Jesus Christ, in the seventh chapter of Daniel. For they are of the same essence. And hereby is signified wisdom and age, and also the divinity and deity of Christ. And because Christ is God, therefore he is the head of the church, ministering to the body spiritual wisdom, and all celestial gifts. Ephesians 5. The Pope of Rome, that most wicked man of sin, what a head is he then? Without life, without brains, utterly foolish. As he is described in the eleventh chapter of Zechariah. It is a shame that we will not see these things, being ever blind. Christ is everlasting, omnipotent, and knows all things. And he may be the health and head of the body. “In the beginning was the Word, and the Word was with God,” says he. Christ himself: “Before Abraham was, I am,” says he. Therefore, the heretics lie, denying Christ to be very God, of the same substance with the Father. He is the dominion of God, all things are subject to him. Ephesians 1. And he himself fulfills all things, just as he is present with his church.

Now his eyes are not darkened nor blind, but keen and bright. For Christ knows all things. Christ’s eyes are watchful; nothing is hidden from him, he sees all things that are done, both good and evil. And he sees in order that he may judge and require. He is light in darkness, and the sight of Christ is to good men joyful in perils. Finally, the judgments of Christ are righteous. The prophet David says, “The eyes of the Lord are upon the just, and his ears are to their prayer.” Again, “The face of the Lord is against them that do evil.” And just as the head is not plucked from the body, so Christ cannot be absent from his church. And considering that his eyes are quicksighted, and that the Lord foresees all our things, and has the charge over us, how is he absent from his church? What need is there of any vicar?

\index{Erasmus of Rotterdam}\index{Ezekiel (Book of)}And the festivals of the Lord are like copper, or like brass and frankincense burning in a furnace. For Chalcolibaum is a word compounded of brass and frankincense. This Erasmus notes, and that Swidas shows also the same, that there is a kind of copper more precious than gold: which he says is made of saltpeter and of a stone. Pliny, in the thirty-fourth and second chapter, calls it a kind of brass, which was dug out of the veins of the earth, and in times past had great price. It seems to me to be the same which in the first and tenth of Ezekiel is called Hasmal, a present remedy against poisons. For if wine intoxicated is put into a cup thereof, it will hiss. And so the death and poison are detected.

The most blatant, fiery spectacle signifies the fellowship and the ways of the Lord, blameless, his judgments right and just. And that he so walks in the church and governs all things, that in the meantime all uncleanness be detected and consumed, but he himself remains always most holy and pure. For fire purges. God is a consuming fire.

But the voice of Christ is as it were the noise of many waters, not so much because all nations and people commend and praise him, but because the Gospel and word of God have come into the whole world. This voice also, mighty kings could less assuage and appease than they could stop the gushing of waters or contain the winds in sacks. Therefore, by the power of preaching, the Lord is always present in his church.

The hand is an instrument of all instruments, especially the right hand. In this, Christ holds seven stars, namely, seven prelates or pastors of churches in Asia. And even all the bishops throughout the whole world, Christ by his power gives to us as pastors, and instructs, comforts, confirms, and defends them, so that they should preach his word. By this, he joins himself to the church; Christ works through them in the church and preserves them.

The same is more vividly expressed in the words that follow. For a sharp, two-edged sword comes out of the Lord’s mouth. This is the word of God, as Jerome declared very well in the sixth chapter to the Ephesians and fourth chapter to the Hebrews. And this word or sword does not hang upon a wall, nor sticks fast in the sheath, nor hangs by the fringe, but comes out of the mouth. He does not say it came forth or it shall come forth, but it comes forth, as the thing that is in continual operation, or perpetual preaching throughout the world. And it is two-edged, sharp and piercing, as in the heart of the godly unto salvation, so in the heart of the wicked to pain and condemnation. And yet at each day comes forth that sword from the mouth of Christ by the mouths of ministers. The word of Christ is indeed condemned by the world, and is called by many a fable, but it is a sword, and that a sword out of Christ’s mouth. All the unfaithful do find and shall find this, however they resist. With this sword Christ kills the wicked. The effect of this sword is greater than was the sword of Alexander, Pompey, Julius Caesar, or Marius, Attila, or [illegible]. Neither does it matter, though the world now acknowledge it not. It shall accomplish their greatest evil in time to come. Doubtless with this spirit of his, the Lord always continues to comfort and govern the church, so that he is never absent from the same.

Finally, the countenance of Christ shines as the sun does in its greatest strength around noon, when it is bright, clear, and pleasant. By the countenance we recognize men chiefly. Therefore, by the countenance we know Christ. The countenance of Christ is light. Christ, therefore, is light. And that truly a divine and eternal light, enlightening all that they may also be made the children of light, and then the faces of the saints may shine in that day as brightly as the sun and as the face of Christ shone. Matthew 13 and 17. And furthermore, he thus communicates this light unto us (1 John and 1 John). How is it to be thought that Christ should be absent from his church? Thou see how he is present.

And so our Lord Christ has revealed himself to us, making himself visible for salvation, and has opened himself wholly, so that we may understand what he does for us and how he is in his church. In these things are all the mysteries of the Gospel comprehended. For what can you say of Christ that you have not here encompassed? Let us therefore remember them and write them in our minds, that we may embrace Christ, king and bishop, and that we never let him depart from our arms. To him be glory.
